\begin{table}[htbp]\centering\small
\caption{Факторный эксперимент Luna: назначения и наблюдаемые исходы.}
\label{tab:scale-quality}
\begin{tabular}{lrrrrr}\toprule
Условие & Назначено & Получено & Оценено & Успехов & Успех (\%)\\\midrule
D & 3000 & 2996 & 2951 & 1611 & 54.59\\
SN & 3000 & 2998 & 2953 & 1604 & 54.32\\
SR & 3000 & 2997 & 2952 & 1581 & 53.56\\
MN & 3000 & 2983 & 2938 & 1471 & 50.07\\
MR & 3000 & 2987 & 2942 & 1514 & 51.46\\
\bottomrule\end{tabular}
\par\smallskip\noindent Оценено: доступные исходы качества после контроля пригодности, а не число разных задач. Процент успеха равен числу успехов, делённому на число оценённых повторов. В каждом условии назначены три повтора на задачу.
\end{table}
